% BioNexus: A Unified Platform for Deterministic Bioinformatics and AI-Assisted Biological Interpretation
% JOSS-style research article draft

\documentclass[11pt,a4paper]{article}
\usepackage[utf8]{inputenc}
\usepackage{microtype}
\usepackage{amsmath,amssymb}
\usepackage{graphicx}
\usepackage{booktabs}
\usepackage{hyperref}
\usepackage[capitalize]{cleveref}
\usepackage{listings}
\usepackage{xcolor}
\usepackage[margin=1in]{geometry}

\hypersetup{colorlinks=true, linkcolor=blue, citecolor=blue, urlcolor=blue}
\lstset{
  basicstyle=\ttfamily\footnotesize,
  breaklines=true,
  frame=single,
  language=Python,
  keywordstyle=\color{blue},
  stringstyle=\color{green!50!black},
  commentstyle=\color{gray},
}

\title{BioNexus: A Unified Platform for Deterministic Bioinformatics \\
       and AI-Assisted Biological Interpretation}
\author{BioNexus Contributors}
\date{\today}

\begin{document}
\maketitle

\begin{abstract}
Modern biology research increasingly relies on two complementary computational
modalities: \emph{deterministic bioinformatics} (sequence alignment, ORF
finding, primer design, CRISPR guide selection, phylogenetic inference ---
computations with single correct answers given identical inputs) and
\emph{AI-assisted interpretation} (literature synthesis, hypothesis ranking,
natural-language explanation --- probabilistic outputs that must be
evidence-grounded and uncertainty-quantified). Existing tools address one
modality or the other, leaving researchers to manually bridge the two. We
present \textbf{BioNexus}, a unified Apache-2.0 platform that pairs
\textbf{BioKit 2.0} (30 deterministic bioinformatics modules organised in 22
subpackages) with \textbf{NEXUS $\Omega$} (an AI-native Biological
Interpretation Intelligence platform with 10 subpackages) via a
\textbf{bridge layer} of 23 adapter programs that enforce a strict
determinism boundary. The platform ships with 275 passing tests, 72\% line
coverage, a unified CLI, REST API, and an evaluation harness supporting
four production LLM providers (OpenAI, Anthropic, GLM, Ollama). We
demonstrate the platform on three canonical biology workflows ---
literature-grounded ORF interpretation, PCR primer validation, and CRISPR
guide evaluation --- and argue that the enforced separation between
deterministic and probabilistic computation is a necessary architectural
invariant for trustworthy AI in biology.
\end{abstract}

\section{Introduction}

The application of large language models (LLMs) to biology has produced
impressive demonstrations but suffers from a recurring failure mode: LLMs
\emph{fabricate biological facts} when asked to perform computations that
have single correct answers. A model asked to compute the GC content of a
sequence, find open reading frames, or design PCR primers may produce
plausible-sounding but numerically wrong output. The standard remedy ---
prompt engineering and chain-of-thought --- does not eliminate the problem;
it merely reduces its frequency.

A growing consensus holds that the correct architectural response is not
to make LLMs better at arithmetic, but to \emph{forbid them from doing it}.
Deterministic bioinformatics computations should be delegated to validated
software, and the LLM should consume the validated output as evidence. This
separation is easy to state and difficult to enforce: it requires a formal
boundary between the two computational modalities, content-addressed
provenance so AI-modified results cannot masquerade as deterministic, and
an evidence ledger that tracks the origin and confidence of every claim.

\textbf{BioNexus} is our implementation of this architecture. It comprises:

\begin{itemize}
  \item \textbf{BioKit 2.0}: 30 deterministic bioinformatics modules
        organised in 22 subpackages, providing sequence manipulation,
        alignment (Needleman-Wunsch, Smith-Waterman, Gotoh), genome
        assembly (De Bruijn, OLC), BLAST parsing, primer design with
        SantaLucia (1998) nearest-neighbour thermodynamics, CRISPR sgRNA
        design, six-frame ORF finding, population genetics statistics
        (Hardy-Weinberg, Weir-Cockerham F\textsubscript{ST}, linkage
        disequilibrium), structural bioinformatics (PDB parsing, RMSD,
        Kabsch superposition), phylogenetics (UPGMA, Robinson-Foulds),
        and a from-scratch machine learning module (KNN, K-Means, PCA ---
        no scikit-learn dependency).
  \item \textbf{NEXUS $\Omega$}: an AI-native interpretation platform with
        a central intelligence kernel (evidence ledger, claim graph,
        uncertainty quantification, governance gates), five agents
        (literature, validation, experiment, workflow, report), RAG
        infrastructure with citation tracking, knowledge source clients
        for PubMed, NCBI, UniProt, PDB, Ensembl, and Gene Ontology,
        and pluggable LLM providers for OpenAI, Anthropic, Google Gemini,
        Zhipu GLM, and local Ollama models.
  \item \textbf{Bridge}: 23 adapter programs that wrap each BioKit module
        as a NEXUS \texttt{BioKitProgram}, automatically registering its
        output as deterministic evidence in the ledger.
\end{itemize}

\section{Architecture}

\Cref{fig:architecture} shows the platform's layering. The user interacts
through a unified CLI, Python SDK, REST API, or web UI. Requests descend
through the NEXUS orchestrator (which selects an agent), through the
intelligence modules (which may invoke the RAG retriever and knowledge
sources), and terminate at the bridge layer when deterministic computation
is required. The bridge dispatches to BioKit 2.0, which produces a
content-addressed \texttt{BioKitOutput} that automatically enters the
evidence ledger as \texttt{EvidenceClass.DETERMINISTIC} --- the only
evidence class whose confidence does not decay over time.

\begin{figure}[t]
\centering
\begin{verbatim}
                User (CLI / SDK / API / Web)
                              |
                +-----------------------------+
                |     NEXUS Orchestrator      |
                |  (5 agents, intelligence,   |
                |   RAG, memory, reports)     |
                +-----------------------------+
                              |
                +-----------------------------+
                |       Bridge Layer          |
                |  (23 BioKitProgram          |
                |   adapters)                 |
                +-----------------------------+
                              |
                +-----------------------------+
                |       BioKit 2.0            |
                |  (30 deterministic          |
                |   bioinformatics modules)   |
                +-----------------------------+
\end{verbatim}
\caption{BioNexus three-layer architecture. The bridge is the only place
where NEXUS learns about BioKit's concrete API, enforcing the Prime
Directive that AI modules may not modify deterministic outputs.}
\label{fig:architecture}
\end{figure}

\subsection{The Prime Directive}

The platform enforces a single architectural invariant, which we call the
\emph{Prime Directive}:

\begin{quote}
\textit{Never replace deterministic biological computation. All biological
calculations remain inside BioKit or other validated scientific software.
NEXUS exists to interpret, validate, explain, reason, retrieve evidence,
generate reports, quantify uncertainty, orchestrate workflows, and assist
researchers.}
\end{quote}

This is enforced at three levels:

\begin{enumerate}
  \item \textbf{Type level}: \texttt{BioKitOutput} is a frozen Pydantic
        model with \texttt{deterministic=True} fixed by construction.
  \item \textbf{Evidence level}: \texttt{BioKitOutput.as\_evidence()}
        returns an \texttt{Evidence} record with
        \texttt{source\_class=DETERMINISTIC}, the only class exempt from
        confidence decay.
  \item \textbf{Engine level}: any attempt to register a deterministic
        claim with an uncertainty envelope raises
        \texttt{DeterministicClaimHasUncertaintyError}; any attempt to
        register an AI claim without one raises
        \texttt{AIClaimMissingUncertaintyError}.
\end{enumerate}

\section{Implementation}

BioNexus is implemented in Python 3.10+ with strict type hints throughout.
The codebase comprises approximately 12,000 lines of Python across 130
source files. Key engineering choices:

\begin{itemize}
  \item \textbf{Pydantic v2} for all data models, providing serialisation,
        validation, and JSON schema generation.
  \item \textbf{Hatchling} as the build backend, producing a single wheel
        that installs all three packages (\texttt{biokit},
        \texttt{nexus}, \texttt{bridge}).
  \item \textbf{Ruff} for linting and formatting with a strict ruleset
        (pycodestyle, Pyflakes, isort, flake8-bugbear, pyupgrade,
        ruff-specific, flake8-simplify, flake8-bandit, Perflint).
  \item \textbf{pytest} with coverage gating at 65\% (limited by optional
        LLM provider SDKs that cannot be tested without credentials).
  \item \textbf{Hypothesis} property-based tests for sequence invariants
        (e.g., GC fraction $\in [0,1]$, complement is involutive).
  \item \textbf{pytest-benchmark} for performance regression tests on
        hot paths (De Bruijn assembly, Needleman-Wunsch, reverse
        complement).
\end{itemize}

\section{Results}

\subsection{Test suite}

The unified test suite comprises 275 tests across three categories:

\begin{table}[h]
\centering
\begin{tabular}{lrr}
\toprule
Category & Tests & Coverage \\
\midrule
BioKit 2.0 unit + integration & 157 & 82\% \\
NEXUS unit (agents, engine, RAG) & 88  & 68\% \\
Bridge + end-to-end integration & 30  & 91\% \\
\midrule
Total & 275 & 72\% \\
\bottomrule
\end{tabular}
\caption{Test suite breakdown.}
\end{table}

\subsection{Benchmark timings}

End-to-end benchmark timings on a single-core Linux VM:

\begin{table}[h]
\centering
\begin{tabular}{lr}
\toprule
Operation & Mean time \\
\midrule
DNA reverse complement (1 kb) & 42 $\mu$s \\
Needleman-Wunsch (100 bp $\times$ 100 bp) & 2.3 ms \\
De Bruijn assembly (1,000 reads) & 40 ms \\
N50 (10,000 contigs) & 0.7 ms \\
Bridge dispatch (gc\_content) & 0.3 ms \\
Bridge dispatch (find\_orfs, 1 kb) & 1.1 ms \\
\bottomrule
\end{tabular}
\caption{Benchmark timings.}
\end{table}

\subsection{Eval harness}

We evaluated the platform on 5 golden questions covering GC content
interpretation, ORF analysis, primer validation, CRISPR guide evaluation,
and alignment interpretation. Each question was run against each
configured LLM provider; providers without credentials were skipped. The
DummyProvider (deterministic mock) achieved a 60\% keyword-match pass
rate, providing a baseline for evaluating real LLM providers in future
work.

\section{Discussion}

The enforced separation between deterministic and probabilistic
computation addresses the most common failure mode of LLM-based biology
tools: fabrication of numerical results. By routing every biological
computation through BioKit and recording its output as
\texttt{EvidenceClass.DETERMINISTIC}, the platform provides a verifiable
audit trail from any AI-generated claim back to the deterministic
computation that supports it.

\subsection{Limitations}

\begin{itemize}
  \item The bridge currently registers 23 BioKit programs; the remaining
        BioKit modules (e.g., visualisation, phylogeny tree building) are
        not yet exposed through the NEXUS facade.
  \item The evidence ledger is in-memory; production deployments require
        a persistent backend.
  \item LLM provider tests are gated behind credentials; the eval harness
        must be run manually with real API keys.
\end{itemize}

\section{Availability}

BioNexus is available at \url{https://github.com/227182038-rgb/BioNexus}
under the Apache License 2.0. Documentation is hosted at
\url{https://227182038-rgb.github.io/BioNexus/}.

\section{Acknowledgements}

BioNexus is built on the work of the open-source bioinformatics and ML
communities. We acknowledge BioPython, the NCBI E-utilities, UniProt, the
Protein Data Bank, Ensembl, the Gene Ontology Consortium, OpenAI,
Anthropic, Google, Zhipu AI, and the maintainers of Pydantic, FastAPI,
Typer, and Rich.

\end{document}
